We have presented Looped World Models, the first application of looped transformer architectures to world modelling. Our approach addresses a central tension in current world models: generating faithful long-horizon simulations demands deep computation, yet deeper models incur prohibitive deployment costs and are susceptible to compounding rollout errors. By iteratively refining latent environment states through a parameter-shared transformer block with stabilised residual dynamics, LoopWM structurally mirrors the recurrence inherent in physical systems while maintaining a compact parameter footprint. Empirically, LoopWM achieve up to 100× parameter efficiency over conventional approaches without sacrificing prediction quality. Theoretically, we show that spectral-norm constraints on state transitions yield provably stable rollouts, providing formal guarantees that are absent in standard autoregressive world models. Furthermore, our adaptive computation mechanism automatically scales the effective depth of the model to match the complexity of each prediction step, allocating more refinement iterations to dynamically challenging transitions and fewer to predictable ones. Beyond the specific results reported here, we believe this work identifies iterative latent depth as a new scaling axis for world simulation, one that is orthogonal to the conventional axes of model size and data volume. We hope that this perspective opens new directions for building world models that are simultaneously more capable, more efficient, and more stable over extended horizons.

